\label{chap:83-12}

\section{Le noyau de transduction de Hensel--Witt}
\label{p12-83-c12-s1:chap:nthw}
\index{Noyau de transduction de Hensel--Witt}

\readerquestion{Où les cylindres, la monodromie, les
cinq régions, les hexades, les octades et les lecteurs physiques se croisent-ils réellement ?}

\subsection{Noyau de transduction de Hensel–Witt : microrégions de \(\pi\), holonomie \(\Gamma_9\), registre dodécaphasique et transport hexade–octade}

Le NTHW est le paquet de données et d’applications
\[
\NTHW=
\bigl(
\mathscr R_\pi^{\rm micro},
\Gnine,
\Cdoce,
K,
\mathcal W_{12\to24},
\mathcal I_{6\to8}
\bigr),
\]
avec les compatibilités suivantes :
\begin{enumerate}
  \item les régions de Hensel partagent un mot quotient et conservent
  des mémoires distinctes ;
  \item \(\Gnine\) transporte l’orientation et produit une monodromie ;
  \item \(\Cdoce\) organise l’observable signé et le sceau \(K\) ;
  \item la lecture de Witt transforme les supports du même bord en hexades et en
  étoiles ;
  \item l’extension \(W_{12}\to W_{24}\) produit l’interface
  hexade--octade ;
  \item les lecteurs archimédien, exceptionnel et dimensionnel reçoivent
  des invariants du même état.
\end{enumerate}

\begin{figure}[p]
\centering
\resizebox{\textwidth}{!}{\input{figures/fig_nucleo_hensel_witt}}
\caption{Le NTHW comme interface typée. La tâche mathématique centrale consiste à faire
commuter le diagramme, et non à accumuler des coïncidences autour d’un nom.}
\label{p12-83-c12-s1:fig:nthw}
\end{figure}

Le NTHW donne un nom formel au lieu auparavant décrit comme « pierre de
Rosette ». La substitution est conceptuelle : une pierre traduit par ressemblance ; une
transduction précise l’entrée, la sortie, la donnée conservée et la perte
d’information.

\subsection{Du code ternaire au système de Witt}

La frontière réversible d’APP sélectionne, avec des jauges déclarées, une matrice
\(A_W\) sur \(\FF_3\) qui satisfait
\[
A_WA_W^T=2I\pmod3.
\]
Le code graphe
\[
C_W=\{(q,qA_W):q\in\FF_3^6\}
\]
est le code de Golay ternaire étendu. Son énumérateur de poids est
\[
1+264y^6+440y^9+24y^{12}.
\]
Les \(264\) mots de poids six, identifiés au signe près, donnent \(132\)
supports. Chaque sous-ensemble de cinq points appartient à une unique hexade :
\[
W_{12}=S(5,6,12).
\]

Le champ de \(495\) involutions étoile engendre un groupe d’ordre
\(95\,040\), reconnu comme \(M_{12}\). Une étoile isolée n’engendre pas le
groupe. Le repère HMT ordonné possède un stabilisateur trivial et fixe donc une
jauge dans l’action.

\subsection{Hexade, octade et corridor \texorpdfstring{\(90/120\)}{90/120}}

Soit \(H\) une hexade et \(O_H\) une octade qui la contient. En marquant un point
et une paire de points de l’hexade, on obtient
\[
F_6(H)=H\times\binom H2,
\qquad
|F_6|=6\binom62=90.
\]
Si le point marqué peut parcourir l’octade :
\[
F_8(H)=O_H\times\binom H2,
\qquad
|F_8|=8\binom62=120.
\]
L’inclusion \(H\hookrightarrow O_H\) induit \(F_6\hookrightarrow F_8\).
Son défaut normalisé est
\[
\boxed{
\frac{90-120}{24\binom62}
=\frac{6-8}{24}
=-\frac1{12}.}
\]
Ce théorème combinatoire précède la lecture angulaire de Barbero. Les
deux objets partagent l’interface \(6\to8\), mais ne sont pas identiques.

\subsection{\(4\) rutas hacia -1/12}

La construction contient quatre réalisations qui doivent être exposées ensemble et non
comptées comme quatre preuves indépendantes d’une même provenance.

\subsubsection{Opérateur élémentaire du cycle dodécaphasique et mise à jour de la phase signée}

Une dent de \(\Cdoce\) mesure \(30^\circ\). Alors
\[
\frac{30}{120}-\frac{30}{90}
=\frac14-\frac13
=-\frac1{12}.
\]
C’est l’ombre angulaire du motif \(90/120\).

\subsubsection{Matrice de Gram des sections survivantes et critère de positivité}

Dans le saut certifié \(60\to81\), deux fibres ont pour tailles \(0\) et
\(12\). L’opérateur de Gram est
\[
K_{60,81}=\operatorname{diag}(0,12).
\]
Sur le sous-espace vivant, sa pseudo-inverse vaut \(1/12\) ; avec l’orientation
négative de bord,
\[
E_{\rm surv}=-K_{60,81}^{+}
\]
produit \(-1/12\).

\subsubsection{Extracteur d'échelle triadique}

La réalisation par un extracteur d’échelle triadique est la construction par
seconde différence développée intégralement dans
\cref{sec:c12-segunda-diferencia-triadica}. Cette première apparition fixe sa
fonction dans l’inventaire des réalisations ; la résidence indiquée conserve
les définitions de \(T_L\), \(a(L)\), \(A_1(L)\), \(C(L)\), l’indépendance
du résidu à l’égard de \(L\) et le statut exact de la normalisation.

\subsubsection{Corridor modulaire entre cartes dodécaphasiques et conservation de l'invariant de phase}

La fonction êta satisfait
\[
\frac{\eta(\tau)}{\eta(3\tau)}
=q^{-1/12}\prod_{n\ge1}(1+q^n+q^{2n})^{-1}.
\]
L’exposant provient de \(1/24-3/24=-1/12\). En prenant douze copies :
\[
t_3(q)=\left(\frac{\eta(\tau)}{\eta(3\tau)}\right)^{12}
=q^{-1}-12+54q-76q^2-243q^3+\cdots.
\]
Le coefficient \(54\) coïncide avec le défaut APP. Le corridor
\[
j=\frac{(t_3+27)(t_3+243)^3}{t_3^3}
\]
se relie ensuite à la branche modulaire classique.

\begin{exactbox}[title={Ce qui a été démontré dans le corridor}]
Le défaut APP \(54\), l'exposant êta \(-1/12\), le défaut
hexade--octade et les identités modulaires une fois \(t_3\) fixé sont exacts. L'affirmation
forte est qu'ils sont tous des images naturelles d'un unique opérateur
TPK. Cette naturalité doit être démontrée au moyen du diagramme du NTHW ; elle ne se
déduit pas de la seule répétition des mêmes entiers.
\end{exactbox}

\subsection{Défaut des puissances de nombres premiers et opérateur nombre}

Pour une puissance d'un nombre premier \(p^m\), le défaut complet d'APP se généralise en
\[
\Delta_{p^m}^{\rm full}
=\frac{m(p-1)}2p^{2m-1}.
\]
En particulier,
\[
\Delta_3=3,
\qquad
\Delta_9=54,
\qquad
\Delta_{27}=729.
\]
En normalisant,
\[
a_{p,m}
=\frac{2\Delta_{p^m}^{\rm full}}{(p-1)p^{2m-1}}
=m.
\]
Avec \(x=p^{-s}\) :
\[
\sum_{m\ge1}a_{p,m}x^m
=\sum_{m\ge1}mx^m
=\frac{x}{(1-x)^2}
=x\frac{\dd}{\dd x}\frac1{1-x}.
\]
Le facteur \((1-x)^{-1}\) est le facteur local d'Euler ; la dérivée insère
l'opérateur nombre. Un pont exact apparaît ainsi :
\[
\boxed{
\text{tour de défauts APP dans }p^m
\longleftrightarrow
\text{insertion de l'opérateur nombre dans le facteur local d'Euler}.}
\]

\subsection{Conditions de naturalité conjointe pour la transduction archimédienne, exceptionnelle et dimensionnelle}

Le NTHW sera un théorème complet de transduction lorsqu'il existera une famille
d'invariants \(\mathcal I\) et des lecteurs
\[
\mathfrak R,\quad\mathfrak E,\quad\mathfrak D
\]
tels que
\[
\mathcal I
=\mathcal I_{\rm arq}\circ\mathfrak R
=\mathcal I_{\rm exc}\circ\mathfrak E
=\mathcal I_{\rm dim}\circ\mathfrak D,
\]
et lorsque les applications de région, de sceau, d'incidence et d'action seront dérivées du même
transport. Les preuves finies fournissent déjà des candidats : monodromie,
support \(B_K\), motif d'emprunts, parité two-way et dénombrements \(90/120\).

\begin{programbox}[title={Preuve exécutable 108--144--90--120}]
Étiqueter chacune des \(1008\) routes de \(\pi\) par orbite CT108 et par
complétion de Witt ; calculer les fibres de l'application ; vérifier si les quatre
régions de multiplicité \(144\) réalisent le facteur \(8/6\) sur des blocs
de \(108\) ; et vérifier si la région de \(432\) transporte la monodromie. Un
seul contre-exemple à la compatibilité des étiquettes réfuterait ce pont dans
sa forme actuelle.
\end{programbox}

\begin{takeawaybox}
Le NTHW est le véritable centre du récit parce que le même bord peut s'y lire
comme cylindre, monodromie, sceau, incidence et action. Sa puissance ne tient pas
à la ressemblance de nombreux nombres, mais au fait que plusieurs applications finies
partagent déjà des supports et des défauts exacts. Le travail suivant consiste à prouver la
naturalité conjointe de ces cartes.
\end{takeawaybox}


\section{Inclusion marquée de l'incidence hexade--octade et préservation de ses multiplicités}

La dodécade et l’alignement du
\cref{p12-83-c17-s1:sec:w12-w24-elevacion} étant fixés, soit \(H\) une hexade dans les coordonnées HMT,
soit \(\ell(H)\subset D\) son image et soit
\(O_H=\iota_D(\ell(H))\) son unique relèvement. En marquant un point
et une paire de points de l’hexade résiduelle, on obtient
\[
F_6(H)=\ell(H)\times\binom{\ell(H)}2,
\qquad
|F_6|=6\binom62=90.
\]
Si le point marqué peut parcourir l’octade :
\[
F_8(H)=O_H\times\binom{\ell(H)}2,
\qquad
|F_8|=8\binom62=120.
\]
L’inclusion \(\ell(H)\hookrightarrow O_H\) induit le morphisme d’incidence
\[
 \mathcal I_{6\to8}:F_6(H)\hookrightarrow F_8(H).
\]
En normalisant la différence par le facteur
choisi \(24\binom62\), on obtient
\[
\boxed{
\frac{90-120}{24\binom62}
=\frac{6-8}{24}
=-\frac1{12}.}
\]
L’égalité est un défaut normalisé de l’interface combinatoire
\(6\to8\). Le facteur \(24\binom62\) fait partie de la définition de cette
normalisation ; l’inclusion à elle seule détermine la différence \(-30\).
L’opérateur \(\mathcal I_{6\to8}\) a pour domaine les configurations
marquées d’incidence ; ce n’est ni l’extracteur transversal de l’état
dodécaphasique ni la famille des moments angulaires.

\section{Classification par domaine des \(5\) réalisations algébriques de \(-1/12\)}

Outre le défaut normalisé de l’inclusion hexade--octade, apparaissent
cinq réalisations opératorielles du même rationnel. Leurs domaines sont
précisés séparément ; l’égalité des valeurs n’identifie pas les
opérateurs sans un entrelaceur supplémentaire.

\subsection{Différence de périodicités dodécaphasiques}

Une dent de \(\Cdoce\) mesure \(30^\circ\). Alors
\[
\frac{30}{120}-\frac{30}{90}
=\frac14-\frac13
=-\frac1{12}.
\]
C’est la différence orientée entre les périodes de \(90^\circ\) et
\(120^\circ\) dans cette réalisation.

\subsection{Pseudo-inverse de l’opérateur de Gram sur la composante vivante}

Dans le saut exact \(60\to81\), deux fibres ont pour tailles \(0\) et
\(12\). L’opérateur de Gram est
\[
K_{60,81}=\operatorname{diag}(0,12).
\]
Sur la composante non nulle de son image, la pseudo-inverse vaut \(1/12\) ; avec l’orientation
négative de bord,
\[
E_{\rm surv}=-K_{60,81}^{+}
\]
produit \(-1/12\).

\subsection{2e différence entre échelles triadiques}
\label{sec:c12-segunda-diferencia-triadica}

Définissons
\[
T_L=\sum_{n=1}^{L-1}n(L-n)=\frac{L(L^2-1)}6,
\qquad
a(L)=-\frac{T_L}{L^2}=-\frac L6+\frac1{6L}.
\]
La première soustraction d’échelle est
\[
A_1(L)=a(L)-\frac{a(3L)}3=\frac4{27L},
\]
et la seconde
\[
C(L)=LA_1(L)-\frac{3L}{3}A_1(3L)=\frac8{81}.
\]
Le résidu \(8/81\) est indépendant de \(L\). Pour obtenir
\(-1/12\), on applique la normalisation déclarée
\[
R(L)=-\frac{27}{32}C(L)=-\frac1{12}.
\]
Le calcul dérive \(8/81\) ; il ne dérive pas à lui seul le facteur \(-27/32\).

\subsection{Composition de transformations modulaires et préservation de l'état signé}

La fonction êta satisfait
\[
\frac{\eta(\tau)}{\eta(3\tau)}
=q^{-1/12}\prod_{n\ge1}(1+q^n+q^{2n})^{-1}.
\]
L’exposant provient de \(1/24-3/24=-1/12\). En prenant douze copies :
\[
t_3(q)=\left(\frac{\eta(\tau)}{\eta(3\tau)}\right)^{12}
=q^{-1}-12+54q-76q^2-243q^3+\cdots.
\]
Le coefficient \(54\) coïncide avec le défaut APP\@. L’identité rationnelle
\[
j=\frac{(t_3+27)(t_3+243)^3}{t_3^3}
\]
établit la connexion avec la branche modulaire classique.

\subsection{Défaut des projecteurs de différence cyclique}

Soit \(S\) le décalage cyclique sur \(\QQ^{12}\) et définissons
\[
 D_r=I-S^r.
\]
Le noyau de \(D_r\) est formé des vecteurs constants sur chaque orbite
de \(S^r\) ; par conséquent,
\[
 \dim\ker D_r=\gcd(12,r).
\]
Soient \(\Pi_{\ker D_3}\) et \(\Pi_{\ker D_4}\) les projecteurs rationnels sur
\(\ker D_3\) et \(\ker D_4\), et soit
\(\tau(T)=\operatorname{Tr}(T)/12\) la trace normalisée. Alors
\[
 \boxed{
 \tau(\Pi_{\ker D_3}-\Pi_{\ker D_4})
 =\frac{\operatorname{rank}\Pi_{\ker D_3}
       -\operatorname{rank}\Pi_{\ker D_4}}{12}
 =\frac{3-4}{12}
 =-\frac1{12}.}
\]
Cette réalisation exprime le rationnel comme défaut transversal des pas
trois et quatre dans le cycle à douze phases.

\begin{remark}[Domaines des réalisations]
Le défaut APP \(54\), l’exposant êta \(-1/12\), le défaut
combinatoire du relèvement hexade--octade et les
identités modulaires une fois \(t_3\) fixé sont exacts. La valeur
de la pseudo-inverse, l’extracteur normalisé et le défaut des projecteurs sont également exacts.
Chaque égalité conserve son domaine et son application correspondante.
\end{remark}
